\chapter{Les modèles, et comment on les choisit}
\label{ch:modeles}

\section{Les modèles retenus}

\begin{center}\small
\begin{tabularx}{\textwidth}{@{}L{2.3cm} L{4.3cm} l l Y@{}}
\toprule
\textbf{Rôle} & \textbf{Modèle} & \textbf{Paramètres} & \textbf{Taille} & \textbf{Où il tourne} \\
\midrule
Filtre d'entrée (étage 1) & Lexique de mots-clés par famille & --- & quelques Ko & CPU \\
Encodeur (étage 1) & \f{multilingual-e5-small} (intfloat) & 118 M & $\approx$\,0,47 Go & CPU ou GPU \\
Tête de classification & Régression logistique & $\approx$\,38\,000 & $<$ 1 Mo & CPU \\
Base de comparaison & TF-IDF + régression logistique & --- & quelques Mo & CPU \\
Formulation (étage 5) & \f{ministral-3:3b}, Q4\_K\_M & $\approx$\,3 G & 3,0 Go & GPU (RTX 3060) \\
Escalade & \f{ministral-3:14b}, Q4\_K\_M & $\approx$\,14 G & 9,1 Go & GPU, à la demande \\
Appétence (étage 3) & Régression logistique & $\approx$\,50 & $<$ 1 Mo & CPU \\
Préparation, hors ligne & \f{ministral-3:14b} ou plus gros, via vLLM & --- & --- & DeepEirb, une fois \\
\bottomrule
\end{tabularx}
\end{center}

La tête de classification compte environ 38\,000 poids : 99 classes (14 familles et 85
intentions) fois 385 (les 384 dimensions du vecteur, plus un biais). Les tailles des modèles
Ollama sont celles de la bibliothèque Ollama ; celles des modèles Ministral 3 incluent un encodeur
d'images que nous n'utilisons pas.

\begin{attention}
\textbf{La mémoire de la RTX 3060 est de 12 Go.} \f{ministral-3:3b} (3,0 Go) et
\f{ministral-3:14b} (9,1 Go) ne tiennent pas ensemble avec leurs caches de contexte. Trois options,
par ordre de préférence :
\begin{enumerate}
  \item l'escalade est envoyée au serveur vLLM de DeepEirb, qui sert le gros modèle à toutes les
        équipes ; le poste ne garde que le petit ;
  \item l'escalade utilise \f{ministral-3:8b} (6,0 Go), qui tient à côté du 3B ;
  \item un seul modèle est chargé à la fois ; l'escalade, rare, accepte quelques secondes de
        rechargement.
\end{enumerate}
Dans tous les cas, le contexte est réglé au plus juste (\f{num\_ctx} à 4\,096 au lieu de 16\,384) :
moins de mémoire et des réponses plus rapides, comme l'a rappelé Denis Barthou.
\end{attention}

\section{\og Le plus petit modèle qui suffit \fg, rendu opérationnel}

La règle commune ne dit pas comment on sait qu'un modèle \og suffit \fg. Nous le décidons avant
de mesurer, pour ne pas choisir après coup le résultat qui nous arrange. La démarche est celle
qu'un organisateur nous a conseillée : partir d'un modèle qui marche, puis descendre.

\begin{cle}
\textbf{Protocole.} On commence par un modèle qui fait clairement l'affaire : il fixe la référence
de qualité. On descend ensuite en taille, candidat par candidat, sur le même jeu et avec le même
seuil fixé à l'avance. \textbf{On retient le plus petit qui tient encore le seuil.} Un modèle plus
gros n'est retenu que si tous les plus petits échouent --- c'est la définition exacte de \og un
modèle plus gros seulement si le gain est mesuré \fg. Toute l'échelle est évaluée : la courbe
qualité, taille et énergie qui en sort va dans le support.
\end{cle}

\paragraph{Pour l'encodeur.} Candidats : \f{paraphrase-multilingual-MiniLM-L12-v2} et
\f{multilingual-e5-small} (118 M chacun), puis \f{multilingual-e5-base} (278 M). Seuil : rester à
moins d'un point de F1 macro du meilleur candidat testé, au niveau des familles, sur le test
français.

\paragraph{Pour la formulation.} Référence : \f{ministral-3:14b} (9,1 Go, déjà installé), ou un
modèle d'environ 30B sur DeepEirb si Denis Barthou en sert un. Puis, en descendant :
\f{gemma4:e2b} (7,2 Go, déjà installé), \f{ministral-3:8b} (6,0 Go), \f{ministral-3:3b}
(3,0 Go), \f{qwen3:1.7b} (1,4 Go, réflexion désactivée), \f{gemma3:1b} (0,8 Go),
\f{qwen3:0.6b} (0,5 Go) et \f{gemma3:270m} (0,3 Go). Jeu : 30 situations sous forme de faits structurés, tirées des parcours. Seuil :
\textbf{zéro nombre inventé après garde-fou} et au moins 27 réponses sur 30 jugées correctes et
claires en français par deux membres de l'équipe. Notre pari est que \f{ministral-3:3b} sera le plus
petit à tenir le seuil ; la mesure peut descendre plus bas.

\paragraph{Les réglages de génération.} Température basse (0,2), au plus 200 tokens produits,
réflexion désactivée sur les modèles qui en ont une, modèle gardé en mémoire pendant la démo
(\f{keep\_alive}) pour ne pas payer le chargement devant le jury.

\section{Le gros modèle, à l'instant précis}

CGI a encouragé l'usage d'un gros modèle \og à un instant T, pour un besoin spécifique \fg. Nous en
faisons deux usages, et seulement deux.

\paragraph{Hors ligne, une fois : préparer les données.} Traduire les 13\,083 phrases de Banking77
et générer les exemples des nouvelles intentions est un travail de qualité, fait une seule fois.
C'est le bon endroit pour un gros modèle : son coût est payé une fois, mesuré avec CodeCarbon, et
\emph{amorti} sur toutes les requêtes futures. Le pitch donne les deux chiffres : l'énergie totale
de la préparation, et l'énergie par requête en service.

\paragraph{En ligne, rarement : lever un doute.} Quand le classifieur reste sous le seuil après une
question de clarification. Le taux d'escalade est mesuré sur les scénarios et affiché ; une
escalade fréquente signalerait un classifieur insuffisant, pas un besoin de gros modèle.

\section{La traduction de Banking77}

\paragraph{La méthode.} Traduction par lots avec une consigne et un glossaire : adapter au contexte
bancaire français plutôt que traduire mot à mot.

\begin{center}\small
\begin{tabular}{@{}l l@{}}
\toprule
\textbf{Anglais (Banking77)} & \textbf{Français attendu} \\
\midrule
\emph{top up} & alimenter son compte, recharger \\
\emph{transfer} & virement \\
\emph{direct debit} & prélèvement \\
\emph{statement} & relevé de compte \\
\emph{PIN} & code confidentiel \\
\emph{ATM} & distributeur \\
\emph{card payment declined} & paiement par carte refusé \\
\emph{exchange rate} & taux de change \\
\bottomrule
\end{tabular}
\end{center}

\paragraph{Le contrôle.} Contrôles automatiques sur tout le jeu (la sortie est bien en français,
le rapport de longueur avec l'original est plausible, pas de ligne vide) ; relecture humaine de
100 phrases tirées au hasard, avec le taux de phrases acceptables reporté dans la fiche de données.

\paragraph{Le repli.} Si DeepEirb n'est pas accessible lundi soir, la traduction passe par
\f{Helsinki-NLP/opus-mt-en-fr}, un modèle de traduction dédié d'environ 75 millions de paramètres,
très rapide sur la RTX 3060. Il connaît moins bien le jargon bancaire : le glossaire est alors
appliqué en post-traitement, et la relecture porte sur le jeu de test en priorité.

\section{Ce qu'il faut demander lundi soir}

\begin{itemize}
  \item \textbf{À Denis Barthou} : l'accès SSH à DeepEirb, et un modèle servi par vLLM pour la
        préparation et l'escalade (\f{ministral-3:14b} au minimum, un modèle plus gros si
        disponible pour la traduction).
  \item \textbf{Sur les postes, en filaire} : télécharger \f{ministral-3:3b}, \f{qwen3:1.7b},
        \f{gemma3:1b} et \f{multilingual-e5-small}. Les slides demandent d'éviter les gros
        téléchargements en wifi.
  \item \textbf{Aux organisateurs} : la version enrichie de Bank Marketing, et la confirmation de
        la machine de démonstration (DeepEirb ou poste de la salle).
\end{itemize}
